\section{Appendix: Liturgical Resolution}
This is the U.S. Thirty-second Sunday in Ordinary Time, Year A, 8 November 2026, Lectionary no. 154 (PC-S58-A): an ordinary Sunday in green, with Gloria and Creed. The governing books are the U.S. \textit{Roman Missal, Third Edition} (2011), based on the Latin third typical edition's emended reprint (2008), and the U.S. \textit{Lectionary for Mass}, second typical edition, Volume I. The national occurrence is established; no particular church's calendar or local choice has been supplied.

The two second-reading lengths are appointed alternatives, as are the two Communion antiphons. Neither is asserted as locally selected or assigned exclusively to Year A. Permitted chant replacements, sprinkling, preparation chant, Preface and Eucharistic Prayer, and concluding blessing remain unselected. Eucharistic Prayer IV retains its integral Preface. No Sequence, short Gospel, proper Offertory antiphon, Sunday-specific Eucharistic Prayer insert or special ritual substitution is established. The expansive study's branch table gives the detailed rubric locators.

The common Week 32 prayers serve all three Sunday cycles. The \textit{General Introduction to the Lectionary}, nos. 103--107, describes the Ordinary Time Old Testament selection in relation to the Gospel and the semi-continuous apostolic and Gospel courses. The complete-Mass interpretations here are source-grounded theological synthesis; each Father is cited for his own passage.

\section{Appendix: Scope and Qualifications}
This concise study derives from the expansive study \textit{Ready to meet the Bridegroom} and its research dossier under the same canonical owner. It adds no research claim. The source edition is rights-limited: full approved English and modern Latin orations, the approved response and acclamation, and other protected liturgical bodies are omitted. Incipits identify prayers; explanations are commentary, not substitute translations. Biblical wording used for study is Douay--Rheims, Challoner revision, Gutenberg 1581, a composite electronic witness not identified as the 1899 American edition. It is public-domain study English, not the English proclaimed at Mass. The approved U.S. Lectionary controls proclamation; a biblical source citation for an adapted antiphon does not reproduce the adaptation.

Authentication of the complete 2008 altar book remains incomplete. The common source owner records the secondary Latin reproduction, official ICEL/FDLC excerpts, and a U.S. 2011 digital-book candidate without independently authenticated publisher transmission. Analysis of the offerings relies on the inspected Latin and that qualified candidate. The older Collect conclusion in retained excerpts is not current recitation text. Precise source locators and full evidence limits remain in the expansive study's References and the common Week 32 owner.

Gregory's identified Latin digital text was read without facsimile collation. Bellarmine is used through O'Sullivan's historical English abridgment, including its omissions. The research found no checked direct Father or saint exposition of the modern offerings and after-Communion wording; their treatment is original analysis. The oil is not assigned uniquely to one sacrament, and no particular person's final fate is pronounced. This is an aid to Catholic study, not an official liturgical book or an assertion of ecclesiastical approval.

\section{References}
\begin{enumerate}
\item \textit{Douay--Rheims Bible}, Challoner revision, Gutenberg 1581: Wisdom 6:13--17 (appointment 6:12--16); Psalms 22, 62 and 87 Vulgate; Matthew 24--25; Luke 24; 1 Thessalonians 4. The retained Thessalonians witness numbers the long appointment 4:12--17 and the short 4:12--13.
\item USCCB, \href{https://bible.usccb.org/bible/readings/110826.cfm}{Thirty-second Sunday in Ordinary Time, 8 November 2026}, Lectionary 154; \href{https://www.usccb.org/resources/2026cal.pdf}{2026 Liturgical Calendar}, printed pp. 5, 43. \textit{Ordo lectionum Missae} (1981), no. 154, p. 85; \textit{General Introduction to the Lectionary}, nos. 103--107. Appointments, national occurrence and Ordinary Time relationships.
\item Common Missal locators: \href{https://media.musicasacra.com/books/latin_missal2002.pdf}{Latin digital reproduction}, Week XXXII, physical pp. 294--295; \href{https://www.liturgyoffice.org.uk/Missal/Music/Antiphonary.pdf}{ICEL Antiphonary}, p. 82; FDLC, \textit{Mystagogical Reflections}, Ordinary Time 31--Christ the King, pp. 8--9; \href{https://archive.org/download/Roman-Missal-Third-Edition/Roman\%20Missal.pdf}{U.S. 2011 digital-book candidate}, p. 492. The source qualifications above apply; citation grants no redistribution right.
\item Augustine, Sermon 93, §§1--17 (Benedictine numbering; NPNF Sermon 43), \textit{Nicene and Post-Nicene Fathers}, first series, vol. 6; \textit{Expositions on the Psalms}, 22.1--6, 62.1--12 and 87.1--4, NPNF1, vol. 8 (headings XXIII, LXIII, LXXXVIII). Registered CCEL historical English witnesses.
\item John Chrysostom, \textit{Homilies on Matthew}, 77.1--3 and 78.1--2, NPNF1, vol. 10; \textit{Homilies on First Thessalonians}, 6--8, opening pericope expositions only, NPNF1, vol. 13. Registered CCEL historical English witnesses.
\item Gregory the Great, \textit{Homiliae in Evangelia}, 12.1--6 and 23.1--2, \href{https://la.wikisource.org/wiki/Homiliarum_in_Evangelia}{Latin Wikisource witness}; original analytical paraphrase here. The retained export credits Gregorius Magnus and Mizardellorsa and declares CC BY-SA 3.0 for transcription/encoding, distinct from the ancient Latin.
\item Bonaventure, \textit{Commentarius in librum Sapientiae}, chapter 6, Vulgate vv. 13--17, \textit{Opera omnia}, Quaracchi VI (1893), pp. 147--149. Historical facsimile checked by the underlying research.
\item Thomas Aquinas, \textit{Super I ad Thessalonicenses}, chapter 4, lectio 2; \textit{Super Matthaeum}, chapter 25, lectio 1 and opening of lectio 2; identified Latin Corpus Thomisticum witnesses. Modern site bodies are not reproduced.
\item Robert Bellarmine, \textit{A Commentary on the Book of Psalms}, John O'Sullivan's historical English abridgment: Psalm 22 and Psalm 62:1--8.
\end{enumerate}
